\section{损失随规模下降，但只是幂律}
\label{sec:14-Deep-Learning/09-Scaling-and-Adaptation/01}

手上有六个已经训完的模型，参数量从一亿排到七十亿，测试损失从 \(3.4\) 一路降到 \(2.9\)。把这六个点画在对数-对数纸上，它们落在一条相当整齐的直线上，线性拟合的残差小到不值得讨论。会议室里有人据此提出：再放大十倍，损失就能降到能用的水平。

这个推断错在哪里，不在于直线拟合得不好，而在于没有把直线的斜率换算成人能理解的量。整节要做的就是这次换算，以及换算之后应当收回哪些期待。

\subsection{一条经验曲线和它的三个变量}

大量实验反复观察到的形状是：固定其余两个因素、只放大一个因素时，测试损失近似满足

\[
  L(N) \approx L_\infty + \left(\frac{N_c}{N}\right)^{\alpha},
\]

其中 \(N\) 是参数量，\(N_c\) 是一个量纲配平用的常数，\(L_\infty\) 是无论怎样放大都降不下去的那部分损失。把 \(N\) 换成数据量 \(D\) 或计算量 \(C\)，形式不变，只是指数换成 \(\beta\) 或 \(\gamma\)。在语言建模上，这三个指数经验上都落在 \(0.05\) 到 \(0.1\) 这个量级；换成别的模态、别的架构、别的分词方式，数值会变，形状大体保留。

先把话说死：上面这个式子是拟合出来的，不是从任何前提推出来的定理。它没有误差界，没有适用条件的清单，也没有告诉你在什么规模之外会失效。后面会专门讲这一点，但换算要先做，因为换算的结论比来源的问题更刺眼。

\subsection{把指数换算成倍数}

问``损失降一半需要多大规模''，就是解

\[
  \left(\frac{N_c}{N_2}\right)^{\alpha} = \frac{1}{2}\left(\frac{N_c}{N_1}\right)^{\alpha}
  \quad\Longrightarrow\quad
  \frac{N_2}{N_1} = 2^{1/\alpha}.
\]

这个式子里没有任何余地。\(\alpha = 0.1\) 时倍数是 \(2^{10} \approx 10^3\)；\(\alpha = 0.07\) 时是 \(2^{14.3} \approx 2\times 10^{4}\)；\(\alpha = 0.05\) 时是 \(2^{20} \approx 10^{6}\)。也就是说，在这个指数区间里，可减半的损失所对应的规模倍数在一千倍到一百万倍之间，而``再大十倍''能买到的损失下降是

\[
  10^{-\alpha} \approx 0.85 \qquad (\alpha = 0.07),
\]

即百分之十五左右——而且这是在减去 \(L_\infty\) 之后的那一部分上打的折，绝对值只会更小。会议室里那个提议不是过于乐观，是量级上就不成立。

指数小意味着两件相反的事同时为真：规模确实持续起作用，从不饱和；每一份等量的收益都要付出成倍的代价。这两句话必须一起说，只说前一句就成了投资建议，只说后一句就否认了过去十年真实发生过的事。

\subsection{对数纸上的直线会骗人}

\begin{center}
\begin{tikzpicture}[scale=1.0]

% ===== 左：对数-对数坐标 =====
\begin{scope}[shift={(0,0)}]
  \fill[black!10] (0,0.4) rectangle (3.6,3.05);
  \draw[->,>=stealth] (-0.2,0.4) -- (5.2,0.4);
  \draw[->,>=stealth] (-0.2,0.4) -- (-0.2,3.4);
  \node[anchor=north] at (2.4,-0.15) {\footnotesize 规模 \(N\)（对数刻度）};
  \node[anchor=south,rotate=90] at (-1.35,1.9) {\footnotesize 损失 \(L\)（对数刻度）};

  \foreach \i/\lab in {0/{10^{7}},1/{10^{8}},2/{10^{9}},3/{10^{10}},4/{10^{11}},5/{10^{12}}}{
    \draw (\i*0.9,0.4) -- (\i*0.9,0.28);
    \node[anchor=north] at (\i*0.9,0.26) {\scriptsize \(\lab\)};
  }
  \foreach \y/\lab in {2.9/{4.0},1.9/{2.0},0.9/{1.0}}{
    \draw (-0.2,\y) -- (-0.32,\y);
    \node[anchor=east] at (-0.36,\y) {\scriptsize \lab};
  }

  \draw[very thick] (0,2.9) -- (4.5,1.74);
  \fill (0,2.9) circle (1.6pt);
  \fill (4.5,1.74) circle (1.6pt);
  \node[anchor=west] at (0.15,3.15) {\scriptsize 斜率 \(=-\alpha \approx -0.07\)};
  \node[anchor=west] at (2.15,2.62) {\scriptsize \(4.0 \to 1.8\)};

  \draw[dashed] (0,1.9) -- (4.9,1.9);
  \node[anchor=north west] at (0.06,1.86) {\scriptsize 损失减半线};

  \draw[<->,>=stealth] (0,-0.72) -- (0.9,-0.72);
  \node[anchor=west] at (0.95,-0.72) {\scriptsize 往右一格 \(=\) 十倍成本};
\end{scope}

% ===== 右：线性坐标 =====
\begin{scope}[shift={(7.4,0)}]
  \fill[black!10] (0,0.4) rectangle (0.45,3.05);
  \draw[->,>=stealth] (-0.2,0.4) -- (5.2,0.4);
  \draw[->,>=stealth] (-0.2,0.4) -- (-0.2,3.4);
  \node[anchor=north] at (2.4,-0.15) {\footnotesize 规模 \(N\)（线性刻度）};
  \node[anchor=south,rotate=90] at (-1.05,1.9) {\footnotesize 损失 \(L\)};

  \foreach \i in {0,1.5,3,4.5}{
    \draw (\i,0.4) -- (\i,0.28);
  }
  \node[anchor=north] at (0,0.26) {\scriptsize \(0\)};
  \node[anchor=north] at (4.5,0.26) {\scriptsize \(10^{12}\)};
  \foreach \y/\lab in {2.9/{4.0},1.9/{2.6},0.9/{1.2}}{
    \draw (-0.2,\y) -- (-0.32,\y);
    \node[anchor=east] at (-0.36,\y) {\scriptsize \lab};
  }

  \draw[very thick,smooth] plot coordinates {
    (0.00,2.90) (0.01,2.47) (0.05,2.10) (0.20,1.79)
    (0.60,1.60) (1.35,1.46) (2.70,1.38) (4.50,1.32)};
  \fill (0.00,2.90) circle (1.6pt);
  \fill (4.50,1.32) circle (1.6pt);

  \draw[dashed] (0,1.05) -- (5.0,1.05);
  \node[anchor=north east] at (4.98,1.01) {\scriptsize \(L_\infty\)};
  \node[anchor=west,align=left,font=\scriptsize] at (1.30,2.55)
    {从 \(10^{7}\) 到 \(10^{11}\)\\ 全部挤在这一小段里};
  \draw[->,>=stealth] (1.22,2.44) -- (0.50,2.05);
\end{scope}

\end{tikzpicture}
\end{center}

\emph{同一条幂律，左边是对数-对数坐标下的直线，右边是线性坐标下的曲线。左图那条线的斜率就是 \(-\alpha\)，看上去平缓可控；但横轴每往右一格代表十倍成本，五格走完是十万倍，换来的是损失从 \(4.0\) 到 \(1.8\)。右图把同一段历程按真实的资源刻度重画，前四个数量级全部压缩在最左侧那一小段里，其余部分近乎水平。要记住的是：直线的``稳步下降''是坐标系造出来的错觉，读斜率之前先读横轴。}

人对直线有一种近乎本能的信任：直线可以外推，外推可以规划，规划可以写进预算。这套直觉在线性坐标下是对的，在对数坐标下则完全不适用，因为往右挪动同样的距离，付出的代价按几何级数增长。左图上从第四格走到第五格，看起来和从第一格走到第二格是一样长的一步，实际成本相差一千倍。

这个视觉陷阱还有一个更隐蔽的版本。人们常说曲线``还没有饱和''，意思是没有看到明显的拐点。幂律确实永远不饱和——它是一条没有拐点的曲线，任何有限规模处都还在下降。但``不饱和''和``值得继续投''完全是两回事：前者是关于函数形状的陈述，后者要把 \(2^{1/\alpha}\) 这个倍数和你的预算放在一起看。用前者论证后者，是这条曲线最常被误用的方式。

\subsection{它是拟合出来的}

反复强调这条曲线的经验身份，不是学究式的谨慎，而是因为它的身份决定了能拿它做什么。

有几种启发式解释被提出来过。一种把指数和数据分布所在流形的内蕴维数联系起来：如果数据实际集中在一个 \(d\) 维流形附近，用分片近似去覆盖它，误差随样本数下降的速率大致是 \(d\) 的倒数量级，于是 \(d\) 越大指数越小。另一种把训练看成在若干种资源上的联合分配，损失由当下最紧的那一种瓶颈决定，多种瓶颈叠加出接近幂律的整体形状。这些说法能提供直觉，能解释为什么指数普遍偏小，但它们都不是证明：既没有给出 \(\alpha\) 的可计算表达式，也没有划出成立的边界。谁也无法据此预先判断某个新模态上的指数是多少，只能训了才知道。

于是所有基于这条曲线的外推都只能是短距离的。跨一到两个数量级，历史上大体站得住；跨四五个数量级去承诺某个具体损失值，没有任何依据支持，只不过很少有人有机会公开地被证伪。

\subsection{\texorpdfstring{\(L_\infty\) 是信息论意义上的地板}{极限项是信息论意义上的地板}}

式子里 \(L_\infty\) 那一项常被略过，它其实是整条曲线里唯一有原理支撑的部分。

把语言建模的损失写成交叉熵，最小可能值就是数据分布本身的熵。一个理想的、知道真实条件分布 \(P(x_t \given x_{<t})\) 的模型，它的期望损失等于该分布的条件熵，不多也不少——这正是\hyperref[sec:07-Information-Theory/01-Information-and-Entropy/02]{``Shannon 熵与二元熵''}给出的量，而\hyperref[sec:07-Information-Theory/04-Source-Coding/04]{``AEP、典型集与信源编码定理''}说明了它同时也是无损压缩的极限。语言本身是有随机性的：同一个上文之后可以接很多个合理的下文，这份不确定性不是模型的缺陷，是被建模对象的属性。再加上数据里的噪声、抓取错误、重复与标注不一致，这些都被算进 \(L_\infty\)。

所以 \(L_\infty\) 不是``目前还没克服的技术困难''，而是一条地板。算力买的只是 \((N_c/N)^{\alpha}\) 那一项，把它推到零，损失也只降到 \(L_\infty\)。实践中的麻烦是 \(L_\infty\) 无法直接测量，只能和 \(\alpha\)、\(N_c\) 一起从数据里拟合出来，而三个参数同时拟合的结果对拟合区间相当敏感。一份报告里 \(L_\infty\) 的估计值，其不确定性通常比它自己写出来的小数位要大。

\subsection{损失降下去，能力不一定同步上来}

还有一个更麻烦的脱节：损失是连续、平滑地下降的，而人们关心的下游表现常常不是。某些任务的准确率在损失越过某个值之前几乎贴着随机水平，越过之后在不长的区间里明显改善。这个现象被观察到很多次，对它的解释至今没有定论。

一种解释说这反映了模型内部确实发生了某种结构性的变化。另一种解释指出，问题可能出在评价指标上：如果一道题要连对若干步才算得分，那么每一步的正确率平滑上升，整题的得分率也会呈现出被放大的、看起来突兀的跃升——变化的是度量方式，不是模型。把同一批模型换成连续的、部分给分的指标重新评一遍，很多跃升会明显变缓。这条解释解释掉了一部分现象，但没有解释掉全部，两种说法目前都还站着。

对使用者来说，无论哪种解释成立，操作上的结论是一样的：损失是可预测的，能力不是。可以用幂律去估算下一个规模的损失值，不能用它去承诺下一个规模会解锁哪项能力。把这两件事混起来写进立项文档，是这个领域里代价最高的一类误用。

\medskip

到这里，规模与损失的关系已经说清了它能说清的部分：形状是幂律，指数很小，地板由信息论定，能力的关系未定。但真正要做的决策还没被触及——预算是一个数，而 \(N\) 和 \(D\) 是两个数，钱应该怎么在它们之间分？这是一个有约束的优化问题，而且它有干净的解析解。
